% Part: normal-modal-logic
% Chapter: completeness
% Section: completeness-K

\documentclass[../../../include/open-logic-section]{subfiles}

\begin{document}

\olfileid{nml}{com}{cmk}

\olsection{\Log{K} માટે નિર્ધારણ અને પૂર્ણતા}

હવે આપણે કૅનોનિકલ નિદર્શન વડે પૂર્ણતા સ્થાપિત કરવા તૈયાર છીએ.
પૂર્ણતા એ હકીકતમાંથી ફલિત થાય છે કે~$\Sigma$ના કૅનોનિકલ નિદર્શનમાં
સત્ય !!{formula}s ચોક્કસપણે $\Sigma$માંથી !!{derivable} સૂત્રો જ છે.
આ ગુણધર્મવાળા નિદર્શનો $\Sigma$નું \emph{નિર્ધારણ કરે છે} એમ કહેવાય છે.

\begin{defn}
  નિદર્શન $\mModel{M}$ કોઈ સામાન્ય મોડલ તર્ક~$\Sigma$નું
  \emph{નિર્ધારણ કરે છે} જો અને માત્ર જો બધાં !!{formula}s~$!A$
  માટે $\mSat{M}{!A}$ જો અને માત્ર જો $\Sigma \Proves !A$.
\end{defn}

\begin{thm}[નિર્ધારણ]\ollabel{thm:determination}
   $\mSat{M^\Sigma}{!A}$ જો અને માત્ર જો $\Sigma \Proves !A$.
\end{thm}

\begin{proof}
  જો $\mSat{M^\Sigma}{!A}$ હોય, તો દરેક પૂર્ણ
  $\Sigma$-સુસંગત $\Delta$ માટે $\mSat{M^\Sigma}{!A}[\Delta]$.
  તેથી સત્યતા સહાયક પ્રમેય મુજબ દરેક પૂર્ણ $\Sigma$-સુસંગત
  $\Delta$માં $!A \in \Delta$; એટલે
  \olref[lin]{cor:provability-characterization} મુજબ ($\Gamma =
  \emptyset$ લઈને) $\Sigma \Proves !A$.

  ઊલટું, જો $\Sigma \Proves !A$ હોય, તો
  \olref[ccs]{prop:ccs-properties}\olref[ccs]{prop:ccs-closed}
  મુજબ દરેક પૂર્ણ $\Sigma$-સુસંગત $\Delta$માં~$!A$ છે; તેથી
  સત્યતા સહાયક પ્રમેય મુજબ દરેક~$\Delta \in W^\Sigma$ માટે
  $\mSat{M^\Sigma}{!A}[\Delta]$, એટલે કે
  $\mSat{M^\Sigma}{!A}$.
\end{proof}

\Log{K}નું કૅનોનિકલ નિદર્શન~\Log{K}નું નિર્ધારણ કરતું હોવાથી
પરિણામ તરીકે~\Log{K}ની પૂર્ણતા તરત મળે છે:

\begin{cor}\ollabel{cor:Kcomplete}
  પાયાનો મોડલ તર્ક \Log{K} બધાં નિદર્શનોના વર્ગને અનુલક્ષીને પૂર્ણ
  છે, એટલે કે $\Entails !A$ હોય તો $\Log{K} \Proves !A$.
\end{cor}

\begin{proof}
  પ્રતિપક્ષી રીતે, જો $\Log{K} \Proves/ !A$ હોય, તો નિર્ધારણ મુજબ
  $\mSat/{M^{\Log{K}}}{!A}$; તેથી $!A$ માન્ય નથી.
\end{proof}

નિદર્શનોના કોઈ વર્ગને અનુલક્ષીને તંત્ર $\Sigma$ની પૂર્ણતાનો સામાન્ય
કિસ્સો લો, જેમ કે સ્વપ્રતિબિંબિત, સંમિત અને પરંપરિત નિદર્શનોના વર્ગને
અનુલક્ષીને \Log{KTB4}ની પૂર્ણતા. અહીં માત્ર નિર્ધારણ પૂરતું નથી.
તંત્ર $\Sigma$નું કૅનોનિકલ નિદર્શન તે વર્ગનું સભ્ય છે એ પણ બતાવવું
પડશે. કૅનોનિકલ નિદર્શનની રચનામાંથી આ સ્પષ્ટ રીતે ફલિત થતું નથી—અને
તે હંમેશાં સાચું પણ નથી!

\end{document}
